\documentclass[runningheads]{llncs}
\usepackage[T1]{fontenc}
\usepackage{url}
\usepackage{array}
\begin{document}
\emergencystretch=2em
\title{From Detection Scores to Supported Claims: Measurement Audits for
Document Analysis and Retrieval}
\titlerunning{Measurement Audits of Document and Retrieval Evaluation}
\author{Anonymous Authors}
\institute{}
\maketitle
\begin{abstract}
Evaluation scores are used to compare systems, but their interpretation
depends on what an evaluator observes, how observations are scored and
what comparison is supported. We examine these dependencies in three
audits: a staged document evaluator, the released FinRank
evidence-ranking benchmark and researcher-generated runs on two BEIR
datasets. In the staged evaluator, 96 of 118 controlled header changes
were visible. Under a single-responsive-leaf restriction, unscaled
thresholds detected 75/95 versus 20/95 tables, while a post-hoc
scale-matched threshold recovered the same 75/95 tables on saved
outputs. In FinRank, changing the tie policy shifted absolute pairwise
accuracy while preserving the examined hard-versus-random comparison.
For our BEIR runs, paired sign-flip tests with Holm correction supported
two of ten lexical-system comparisons on SciFact and one on NFCorpus. An
exact-tie-aware point-estimate baseline also passed our duplicate and
degradation controls. These cases show how explicit measurement
contracts qualify or preserve particular comparisons.
\keywords{measurement validity, document evaluation, information retrieval,
reproducibility, benchmark auditing}
\end{abstract}
\section{Introduction}\label{introduction}

A lower detection rate can indicate a weaker evaluator, but it can also
arise when the same numerical threshold is applied to scores on
different scales. Likewise, a change in benchmark score need not change
the comparison for which the score was reported; and the highest mean in
a leaderboard does not by itself establish that every apparent ordering
is supported by the data. Such distinctions matter when researchers use
evaluation results to select systems or interpret a benchmark.

We ask: \textbf{which interpretations remain supported when the
observation surface, scoring rule and comparison procedure are made
explicit?} We investigate three evaluation constructions. First,
controlled changes to public financial tables probe a staged document
evaluator. Second, we inspect tie handling and source identity in
FinRank \cite{r1}. Third, we assess researcher-generated lexical
retrieval runs on SciFact and NFCorpus from BEIR \cite{r2}, with a later
dense-model exploration. The BEIR runs are our own constructions, not
the published benchmark leaderboard.

Our contributions are three empirical findings. We explain a
detection-rate difference by a positive score rescaling under a
single-responsive-leaf restriction; show that a particular FinRank
comparison retains its direction despite changes in absolute scores; and
distinguish point-estimate rankings from comparisons supported under a
specified paired statistical procedure.

Prior work has examined metric and construct validity \cite{r3},
\cite{r4}, \cite{r5}. AgentSuite \cite{r6} audits interacting components
in agent benchmarks at broader scale. We complement these perspectives
with artifact-level analyses of observation, score scaling and
comparison in document and retrieval evaluations.

Document analysis evaluation already addresses several related problems.
SCORE \cite{r7} evaluates semantically equivalent but structurally
different generative parsing outputs; Soric et al.~\cite{r8} examine
end-to-end and subtask table-extraction evaluation across heterogeneous
PDFs. Those studies assess document parsing or extraction under their
own evaluation designs. We instead audit the interpretation of
measurements produced by a particular frozen evaluator, including what
its selected mapping surface can observe and what follows from its
scoring scale. Our table inputs are derived from financial QA data, not
page-image OCR outputs. We do not benchmark document extraction models
or claim a superior parsing metric.

The document case is the primary construction: a controlled change to a
table header travels through a mapping rule into a checkpoint score and
then a thresholded decision. Its observation and decision contracts can
be inspected separately. The two external ranking cases complement this
document-centered analysis by examining whether scores and comparison
claims remain stable under other declared conventions; they do not turn
the controlled header study into an external validation of a
document-image analysis method.

\section{Audit Design}\label{audit-design}

We distinguish three questions: whether a change is \textbf{visible} at
a selected evidence surface; whether a specified predicate
\textbf{detects} it; and whether it changes an outcome under a stated
equivalence contract. Visibility, detection and outcome change are
separate observations. We then examine score sensitivity, the stability
of a named comparison and numerical reproducibility.

The designs differ. Case 1 uses controlled perturbations. Case 2
inspects a released benchmark, including source-informed post-hoc
identity analysis. Case 3 uses researcher-generated systems and test
queries from a public benchmark; its original audit choices were not
independently timestamped before the first run. We preserve fixed
outputs and identify later explorations rather than treating them as
prospective confirmation.

The cases target three interpretations. First, a detection-rate
difference does not necessarily imply a ranking difference when the
score scales differ. Second, a shift in absolute accuracy need not
reverse a stated comparison. Third, a higher retrieval mean is a point
estimate; it does not automatically establish every pairwise system
ordering. The following table summarizes what the audit changes in each
interpretation; the cases below specify what is held fixed and what is
varied.

\begin{table}[htbp]
\caption{Interpretations qualified or retained by the three audits.}
\label{tab:interpretations}
\centering
\footnotesize
\renewcommand{\arraystretch}{1.12}
\begin{tabular}{@{}>{\raggedright\arraybackslash}p{.18\columnwidth}>{\raggedright\arraybackslash}p{.34\columnwidth}>{\raggedright\arraybackslash}p{.36\columnwidth}@{}}
\hline
Case & Tempting interpretation & Supported interpretation \\
\hline
Document & Lower detection means weaker ranking & Rates differ; ranks coincide in the single-leaf setting. \\
FinRank & Absolute score changes imply a changed comparison & The random-minus-hard gap stays positive under both inspected tie policies. \\
Retrieval & A higher mean establishes every system ordering & The paired test supports some differences; other pairs are not proven equal. \\
\hline
\end{tabular}
\end{table}

For each construction we identify the input and outcome carrier, the
implementation, the original reporting convention and the comparison of
interest. We use \emph{decision} for an analytical interpretation such
as a comparison's direction; no human-decision intervention was
conducted. The cases use different units and populations, so their
counts are not pooled into an overall audit-success rate.

\section{Controlled Document
Evaluation}\label{controlled-document-evaluation}

Our project-built, frozen staged evaluator has parsing and
question-answering checkpoint trees with twelve scored leaves. Its
seed-42 panel contains 95 TAT-QA-derived source tables \cite{r9} and 118
controlled header changes across six families. Mapping and regression
artifacts form the observed surface. For each injected header, the
visibility analysis identifies its original canonical field: the change
is counted as absorbed when that field still maps to the injected
header, and visible otherwise. This matcher-relative criterion is not
simply a character-level difference. Protocol detection and downstream
harm are separate quantities. The audit holds the evaluated inputs and
scores fixed while comparing declared aggregation and threshold
conventions.

Of 118 injections, 96 were visible. None of the 13 OCR-like changes was
visible on this surface. This observation does not establish 13 harmful
misses: one existing semantic judge classified the cases as
referentially equivalent, without an independent human reference. The
family-level visible/injected counts were 43/47 alias variants, 15/17
date variants, 20/20 mixed-language variants, 0/13 OCR-like variants,
13/16 prefix-noise variants and 5/5 unit-suffix variants. These are
injection counts in the seed-42 panel; the 118 changes are clustered
within 95 tables. They characterize the particular matcher-relative
surface, not generic document defects or actual OCR error rates.

In this construction the QA tree was frozen and only the field-hit-rate
leaf responded. For field hit rate h, the aggregation weights give a
parsing-root score of approximately 6.25 + 3.75h, while the flat drop is
ten times the hit-rate drop. The hierarchical drop is therefore
approximately 0.375 times the flat drop, and their observed rankings
coincide. The original protocols each used a numerical operating
threshold of \textbf{at least 2.0} on their respective scales. At those
thresholds the flat and hierarchical protocols detected 75/95 versus
20/95 tables, or 85/118 versus 13/118 injections. Tables and injections
are distinct units; a rate difference does not by itself establish a
ranking difference.

\begin{table}[htbp]
\caption{Seed-42 original thresholds and post-hoc scale matching; table and injection counts are distinct units.}
\label{tab:original-detection}
\centering
\setlength{\tabcolsep}{4pt}
\begin{tabular}{lrrr}
\hline
Unit & Flat $\ge 2$ & Tree $\ge 2$ & Tree $\ge 0.75$ \\
\hline
Source tables & 75/95 & 20/95 & 75/95 \\
Header injections & 85/118 & 13/118 & 85/118 \\
\hline
\end{tabular}
\end{table}

The flat threshold of 2.0 corresponds to 0.75 on the hierarchical scale.
In a post-hoc check of saved rounded scores, these matched thresholds
select the same 75/95 tables and 85/118 injections as the flat rule. A
separate grid used strict greater-than thresholds of 1.0 and 0.375: in
one recorded fold both selected the same 21 of 24 injections with
experimental H/HI labels. The predicates must not be conflated: 12 of
the 95 tables have a flat drop exactly equal to 2.0.

In this single-leaf setting, a positive rescaling preserves score order;
retaining the same decision rule also requires rescaling its threshold.
The matched-grid example uses saved outputs and does not calibrate a
deployment threshold or establish equivalence with multiple responsive
leaves.

The clean-only contrastive construction compares repeated deterministic
computations on an unchanged record. All 2,382 usable clean records have
zero scores, so every positive threshold has observed clean
false-positive rate zero; this panel does not identify a unique
deployment threshold. Across three seeds, the flat descriptive table-run
ratio was 219/282 using 119 unique source-table identifiers. A post-hoc
UID-cluster bootstrap gives a 95\% interval of {[}0.7117, 0.8364{]}; the
repeated table runs are not 282 independent source tables.

The per-seed contributions to the 219/282 ratio are 75/95, 76/97 and
68/90. The denominators refer to table appearances within each run; the
same source table may contribute to more than one run. The bootstrap
samples source identifiers with replacement, includes all run
appearances for each sampled identifier, and recalculates the pooled
ratio. This uncertainty interval is post-hoc and conditional on the
available source tables and three seeds; it should not be interpreted as
a guarantee for future datasets or seeds. Equal weighting of source
identifiers gives 0.7535 with interval {[}0.6835, 0.8207{]}.

In a separate stage-aligned control, 72 source clusters contributed one
injection each for mapping, retrieval and numeric processing. A
leave-one-cluster-out nearest-centroid classifier on the full response
vector attained macro-F1 and top-1 1.000. The three directly targeted
parsing, retrieval and reconciliation leaves gave the same result;
removing them reduced top-1 to 0.333 under the frozen tie-as-incorrect
rule, with a 0.667 tie rate. The classifier was evaluated by leaving one
source cluster out at a time, so the three injected stages from the
held-out cluster were not included when fitting its centroids. This
ablation tests the source of the observed signal in the controlled
instrument, not error localization on an independent deployed system.
The controlled stage signal thus resides in the directly targeted
leaves. A nested FinQA probe \cite{r10} covered 89 rows from 30 report
clusters; it changed program fields but did not recompute downstream
result fields. Its zero detected and zero changed outcomes cannot
estimate a failure rate for changed answers.

\section{FinRank: Score Sensitivity and Decision
Stability}\label{finrank-score-sensitivity-and-decision-stability}

FinRank ranks evidence from financial filings \cite{r1}. Under
tie-as-positive scoring, a score tie counts for the gold passage; under
strict-win scoring, it does not. The tie comparisons and materiality
threshold were fixed before pair-level outcomes. The following table
shows that random-negative accuracy moves by 1.084 percentage points for
TF-IDF and 1.029 for BM25, exceeding the fixed 1.0-point threshold,
while the random-minus-hard gap remains positive in both systems.

\begin{table}[htbp]
\caption{FinRank pairwise accuracy under two tie policies; gaps are computed before rounding.}
\label{tab:finrank}
\centering
\setlength{\tabcolsep}{6pt}
\begin{tabular}{llrrr}
\hline
System & Policy & Hard (\%) & Random (\%) & Gap (pp) \\
\hline
TF-IDF & Positive & 77.756 & 91.935 & 14.179 \\
TF-IDF & Strict & 77.030 & 90.850 & 13.821 \\
BM25 & Positive & 74.780 & 87.789 & 13.009 \\
BM25 & Strict & 74.094 & 86.760 & 12.666 \\
\hline
\end{tabular}
\end{table}

Gaps are random accuracy minus hard accuracy, computed before display
rounding. The gap's direction and descriptive TF-IDF/BM25 ordering
remain stable under both policies. A source-informed post-hoc analysis
of the 1,185-record release reproduces 442 hard-negative occurrences
across 305 records sharing canonical text with another record's gold
passage. The inspected pooled corpus keeps the first occurrence's
metadata when deduplicating passage text; metadata filtering can
consequently exclude a later record's gold text, as observed for 33
records. These properties were known from the source. In-record and
pairwise metrics reproduce exactly in the audited environment; global
and metadata-filtered metrics differ from the released results by up to
0.19 and 0.16 percentage points, respectively. The cause remains
unresolved.

\section{Retrieval Comparisons}\label{retrieval-comparisons}

We built five lexical configurations on BEIR SciFact and NFCorpus
\cite{r2}: TF-IDF and BM25 parameter pairs (k1, b) of (0.9, 0.4), (1.2,
0.75), (0.5, 0.3) and (1.5, 0.6). We concatenate each document's title
and text, lowercase and tokenize on ASCII alphanumeric runs, use
log-scaled term frequency and smoothed inverse document frequency for
TF-IDF cosine, and resolve equal scores by ascending document ID unless
testing the alternative order. They use the released corpora, qrels and
300/323 judged queries, respectively. The systems are our own runs, and
the two datasets belong to one benchmark family. We compare system
scores on the same queries so that statistical analysis is paired; each
dataset has its own comparison family.

Our initial scoring implementation binarized graded relevance and used a
different MAP denominator. Aligning it with the trec\_eval contract
corrected both definitions. Under graded nDCG@10, TF-IDF moved from
second to third in one NFCorpus ranking; the separate denominator
correction changed TF-IDF's MAP@10 from approximately 0.2200 to 0.1162.
This is a correction of our implementation, not evidence that the
official leaderboard was wrong. For fixed rankings, trec\_eval 10.0-rc3
agreed with 18,690 per-query values within 1e-4: five systems, 623
queries and six nDCG, MAP and recall cutoff combinations. This check
does not validate tie-breaking policy.

The cross-check encodes each fixed rank with a strictly decreasing
score. Its maximum discrepancy was approximately 0.00005, consistent
with the reference program's four-decimal output. It checks metric
arithmetic for the chosen ranking, rather than which document should win
a score tie; the executable was trec\_eval itself, not the pytrec\_eval
Python binding.

For the five systems, 2,000 paired bootstrap resamples produced
intervals excluding zero for 3/10 pairs on SciFact and 4/10 on NFCorpus.
A separate two-sided paired sign-flip procedure with 2,000 Monte Carlo
draws per pair and Holm correction supported 2/10 and 1/10,
respectively. Both sign-flip analyses use the Monte Carlo correction
(b+1)/(M+1), where b is the number of draws at least as extreme as
observed and M the number of draws. Bootstrap and sign-flip are
different inference procedures, so the count difference is not
attributable to Holm correction alone. Inference depends on the relevant
exchangeability or symmetry and query-independence assumptions;
non-significance does not establish equality.

The supported SciFact pairs were BM25(1.5, 0.6) versus TF-IDF (mean
nDCG@10 difference +0.03334; Holm-adjusted p = 0.0100) and BM25(1.2,
0.75) versus TF-IDF (+0.02930; p = 0.0450). On NFCorpus, BM25(1.5, 0.6)
versus BM25(0.5, 0.3) had a difference of +0.00817 (p = 0.0300). A later
200,000-draw sensitivity check on regenerated lexical query scores
retains the same three supported pairs when Holm is recalculated within
the five-system family. It reduces Monte Carlo variability but does not
independently validate the data or inferential assumptions.

Reversing document-ID tie order for BM25(0.9, 0.4) on NFCorpus at k=10
changed the ordered top ten for 36.5\% of queries and top-ten membership
for 25.7\%. Mean nDCG@10 changed from 0.30453945 to 0.30449093,
approximately 0.00005. Query-level changes can therefore be hidden by
the aggregate in this setting; this is not an estimated benchmark-wide
failure rate.

For an intentionally duplicated system and an intentionally degraded
control, an exact-tie-aware point-estimate rule and the paired audit
both passed both controls. The controls do not demonstrate that the
audit outperforms the simpler rule. A subsequent BGE-small-en-v1.5
\cite{r11} exploration gave mean nDCG@10 of 0.7127 versus 0.6622 for
BM25(1.2, 0.75) on SciFact, and 0.3438 versus 0.3062 on NFCorpus. This
is an exploratory retrieval result selected after the original analysis,
not evidence that the audit improved a user decision or that lexical tie
findings transfer to dense models.

The BGE comparison was later included with the five lexical runs in a
separate six-system family of 15 pairs per dataset. With 200,000 Monte
Carlo sign flips per pair and Holm correction within each dataset, the
BGE-minus-BM25(1.2, 0.75) nDCG@10 difference was +0.05049 on SciFact
(unadjusted p = 0.00641; adjusted p = 0.066000) and +0.03759 on NFCorpus
(adjusted p = 0.015540). The SciFact pair crosses the 0.05 threshold
after adjustment within this same test; it is not evidence of equality.
These post-hoc tests are not the original ten-pair family or independent
confirmation. The saved dense outputs contain ranked document IDs rather
than original dense scores, so they cannot establish dense-score tie
frequencies. Numerical reproducibility of the tests does not establish
their inferential assumptions.

The earlier single-pair BGE analysis used normalized embeddings, the
model's query prefix and a 512-token cap. Its uncorrected
paired-bootstrap differences were 0.0505 {[}0.0137, 0.0891{]} on SciFact
and 0.0376 {[}0.0154, 0.0608{]} on NFCorpus. A separate trec\_eval run
agreed on 1,246 dense and BM25 per-query scores to four-decimal print
precision. The later 15-pair test addresses multiplicity, not
prospective confirmation.

\section{Discussion and Conclusion}\label{discussion-and-conclusion}

The three cases show distinct ways in which an evaluation contract
constrains a claim. In the staged evaluator, an apparent rate difference
vanishes after matching thresholds under a restricted score rescaling.
In FinRank, a tested comparison retains direction as absolute scores
change. In our BEIR runs, tie handling and statistical comparison
procedures affect the interpretation of rankings; correction of our own
metric implementation is a separate calibration step.

These findings have limited scope: the first case uses controlled
perturbations to structured tables, without a page-image recognition
experiment or an independent human reference for the OCR-like variants.
FinRank uses lexical baselines and has an unresolved reproduction gap,
and the BEIR analyses are retrospective runs over two datasets with a
post-hoc dense extension. Existing validity frameworks already cover
broader audit principles \cite{r3}, \cite{r4}, \cite{r5}. Our
contribution is a reproducible set of case-specific observations and the
distinction between numerical movement and supported interpretation. We
recommend reporting the measurement object, scoring convention,
comparison and uncertainty together; we do not claim a general audit
guarantee or observed deployment benefit.

\section{Data and Materials}\label{data-and-materials}

The TAT-QA, FinRank and BEIR data used in this study are described in
\cite{r9}, \cite{r1} and \cite{r2}, respectively; the BEIR runs use
released relevance judgments. Study-specific audit scripts,
configuration records and aggregate results are retained as research
artifacts. Access to and redistribution of underlying data remain
subject to the respective source terms.

\section*{Acknowledgment}

Generative AI tools, including DeepSeek \cite{r12} and OpenAI Codex
\cite{r13}, assisted with analysis-code development and drafting and
revision of manuscript text in Sections 1--6. Google Gemini \cite{r14}
provided review suggestions on earlier drafts. The authors are
responsible for the final manuscript and its conclusions.
\begin{thebibliography}{14}
\bibitem{r1} S. Mansouri et al., ``FinRank: An Evidence-Grounded Benchmark for
Financial Question Answering and Retrieval over SEC Filings,''
arXiv:2608.07400, 2026.
\bibitem{r2} N. Thakur et al., ``BEIR: A Heterogeneous Benchmark for Zero-shot
Evaluation of Information Retrieval Models,'' NeurIPS Datasets and
Benchmarks, 2021.
\bibitem{r3} C. Ross et al., ``What makes a good metric? Evaluating automatic metrics
for text-to-image consistency,'' arXiv:2412.13989, 2024.
\bibitem{r4} W. Caban, ``Measurement Without Validity: The Compounding Reliability
Problem in Agentic AI Evaluation,'' arXiv:2608.00794, 2026.
\bibitem{r5} A. M. Bean et al., ``Measuring what Matters: Construct Validity in Large
Language Model Benchmarks,'' arXiv:2511.04703, 2025.
\bibitem{r6} H. Suh et al., ``AgentSuite: Toward More Reliable Agent Evaluation with
a Component-Based Benchmark Auditing Pipeline,'' ICML, PMLR
306:116300--116319, 2026.
\bibitem{r7} R. Li et al., ``SCORE: A Semantic Evaluation Framework for Generative
Document Parsing,'' arXiv:2509.19345, 2025.
\bibitem{r8} M. Soric, C. Gracianne, I. Manolescu, and P. Senellart, ``Benchmarking
Table Extraction from Heterogeneous Scientific PDF Documents,''
arXiv:2511.16134v2, 2026.
\bibitem{r9} F. Zhu et al., ``TAT-QA: A Question Answering Benchmark on a Hybrid of
Tabular and Textual Content in Finance,'' ACL-IJCNLP, 2021, doi:
10.18653/v1/2021.acl-long.254.
\bibitem{r10} Z. Chen et al., ``FinQA: A Dataset of Numerical Reasoning over Financial
Data,'' EMNLP, 2021, doi: 10.18653/v1/2021.emnlp-main.300.
\bibitem{r11} Beijing Academy of Artificial Intelligence, ``bge-small-en-v1.5,'' model
card, revision 5c38ec7c, https://huggingface.co/BAAI/bge-small-en-v1.5.
\bibitem{r12} DeepSeek, ``DeepSeek V4.1,'' https://www.deepseek.com/, accessed
November 2, 2026.
\bibitem{r13} OpenAI, ``Codex,'' https://openai.com/codex/, accessed November 2, 2026.
\bibitem{r14} Google, ``Gemini,'' https://gemini.google.com/, accessed November 2,
2026.
\end{thebibliography}
\end{document}
